%==========================================================================
% CAPÍTULO 2 — REVISÃO DE LITERATURA
%==========================================================================
\chapter{Revisão de Literatura}
\label{chap:revisao}

\section{Visão Computacional e Agricultura de Precisão}

A visão computacional é um ramo da inteligência artificial que visa dotar
sistemas computacionais da capacidade de interpretar e compreender informação
visual proveniente do mundo real. Nas últimas décadas, os avanços em redes
neuronais profundas (\textit{deep neural networks}) transformaram profundamente
o estado da arte nesta área, permitindo alcançar níveis de desempenho
comparáveis ou superiores ao humano em diversas tarefas de reconhecimento e
deteção de objetos~\cite{lecun2015deep}.

No contexto da \textbf{agricultura de precisão}, a visão computacional tem
sido aplicada a um conjunto diversificado de problemas: deteção de pragas e
doenças em folhas~\cite{mohanty2016plant}, estimativa do estado de maturação
de frutos~\cite{sa2016deepfruits}, mapeamento de culturas por drones e
contagem de plantas~\cite{osco2021review}. A viticultura, em particular, tem
atraído crescente atenção da comunidade científica, dado o valor económico da
cultura e a complexidade visual do ambiente de vinha.

\section{Deteção de Objetos com Redes Neuronais Profundas}

\subsection{Evolução dos Modelos de Deteção}

Os modelos modernos de deteção de objetos podem ser classificados em duas
grandes famílias:

\begin{itemize}
  \item \textbf{Modelos em dois estágios} (\textit{two-stage detectors}):
    Como o Faster R-CNN~\cite{ren2015faster}, que primeiro gera propostas de
    regiões e depois as classifica. Tendem a apresentar maior precisão, mas
    com latência mais elevada.

  \item \textbf{Modelos em um estágio} (\textit{single-stage detectors}):
    Como a família YOLO (\textit{You Only Look Once})~\cite{redmon2016you} e
    o SSD~\cite{liu2016ssd}, que realizam a deteção numa única passagem pela
    rede. São mais rápidos e adequados a aplicações em tempo real.
\end{itemize}

Para sistemas embebidos e robóticos, os modelos de um estágio são
preferíveis devido ao menor custo computacional.

\subsection{A Família YOLO e o YOLO26}

A arquitetura YOLO foi introduzida por Redmon \textit{et al.}~\cite{redmon2016you}
e tem sido continuamente aprimorada. O \textbf{YOLOv8}, desenvolvido pela
Ultralytics~\cite{jocher2023ultralytics}, representou um marco importante,
com melhorias na arquitetura do \textit{backbone}, do \textit{neck} e da
cabeça de deteção (\textit{detection head}).

O \textbf{YOLO26}~\cite{hidayatullah2026yolo26}, versão mais recente da
família, introduz avanços significativos sobre o YOLOv8:

\begin{itemize}
  \item Eliminação da \textit{Distribution Focal Loss} (DFL) e adoção de
    inferência \textit{End-to-End NMS-Free}, reduzindo a latência em 43\%
    no modo CPU;
  \item Introdução do \textbf{STAL} (\textit{Small-Target-Aware Label
    Assignment}), que corrige a sub-atribuição de objetos com dimensão
    inferior a 8$\times$8~px — especialmente relevante para a deteção de
    gomos;
  \item Bloco de atenção \textbf{PSABlock} no \textit{neck}, melhorando
    a representação de características sem custo computacional elevado;
  \item Otimizador \textbf{MuSGD}, que combina SGD com atualizações
    estilo Muon, melhorando a convergência.
\end{itemize}

A variante \textbf{YOLO26 Nano} é a mais leve da família, adequada para
execução em tempo real em hardware com recursos limitados, como sistemas
robóticos de campo. O modelo retorna, para cada objeto detetado, uma
\textit{bounding box}, uma classe e um valor de confiança.

\section{Aplicações em Viticultura}

\subsection{Deteção de Uvas e Cachos}

Vários trabalhos têm abordado a deteção de uvas e cachos em imagens de
campo. Santos~\textit{et al.}~\cite{santos2020grape} e Nuske~\textit{et
al.}~\cite{nuske2011estimating} utilizaram visão computacional para estimar
o rendimento das vinhas através da contagem de cachos. A complexidade das
oclusões e a variabilidade de cor entre cultivares são apontadas como os
principais desafios.

\subsection{Deteção de Sarmentos e Estrutura da Planta}

A identificação da estrutura esquelética da videira — tronco, cordões e
sarmentos — é uma etapa anterior à localização dos gomos e tem sido explorada
por Kicherer~\textit{et al.}~\cite{kicherer2015automatic} e Majeed~\textit{et
al.}~\cite{majeed2020deep}. Estes trabalhos demonstram que a segmentação ou
deteção dos elementos estruturais melhora a robustez da localização dos
pontos de corte.

\subsection{Deteção de Gomos}

A deteção específica de gomos (\textit{buds}) é uma tarefa de menor dimensão
aparente, mas de alta dificuldade, dado que os gomos têm dimensão reduzida
e podem apresentar-se parcialmente oclusos por folhas ou outros sarmentos.
Botterill~\textit{et al.}~\cite{botterill2017robot} desenvolveram um sistema
robótico de poda que inclui deteção de gomos com câmeras RGB e de
profundidade. Aquino~\textit{et al.}~\cite{aquino2018vitisberry} focaram-se
em cultivares espanholas e obtiveram resultados promissores com redes
convolucionais.

\section{Plataformas de Anotação de Dados}

A qualidade e consistência das anotações são determinantes para o desempenho
de modelos supervisionados. A plataforma \textbf{Roboflow}~\cite{roboflow2020}
tem emergido como uma das ferramentas mais completas para:

\begin{itemize}
  \item Anotação manual de \textit{bounding boxes} e polígonos;
  \item Aplicação de técnicas de \textit{data augmentation} (rotação, flip,
    ajuste de brilho, corte, mosaico);
  \item Exportação de datasets nos formatos suportados pelos principais
    frameworks (YOLO, COCO, Pascal VOC);
  \item Gestão de versões de dataset e colaboração em equipa.
\end{itemize}

A existência de datasets públicos na plataforma, incluindo conjuntos
provenientes de Espanha e de outros países produtores de vinho, facilita a
aplicação de \textit{transfer learning} e a expansão de dados de treino.

\section{Desafios em Ambientes Vitícolas Reais}

A literatura aponta de forma recorrente um conjunto de desafios específicos
ao ambiente de vinha que dificultam a deteção robusta de objetos:

\begin{description}
  \item[Oclusão e sobreposição:] Os sarmentos de diferentes plantas cruzam-se
    frequentemente, e os gomos podem ficar parcialmente ocultos por outros
    elementos da planta ou por vegetação de fundo.

  \item[Variações de iluminação:] A luz solar direta, a sombra projetada
    pelos sarmentos e as mudanças ao longo do dia afetam significativamente
    a aparência visual dos gomos e sarmentos.

  \item[Fundo complexo:] Em vinhas alinhadas, o fundo de uma imagem de um
    sarmento em primeiro plano inclui tipicamente outros sarmentos, folhas e
    estruturas de suporte, tornando difícil a distinção entre objetos de
    interesse e contexto.

  \item[Variabilidade morfológica:] As dimensões, formas e aspetos dos gomos
    variam consideravelmente entre cultivares, fases de crescimento e
    condições de stresse hídrico.
\end{description}

Estes desafios motivam a necessidade de recolher imagens em condições reais
e de construir datasets com grande diversidade, como descrito no presente
trabalho.
